\chapter{Tabellen}
\label{cha:Tabellen}

Tabellen sind neben Abbildungen ein zweckmäßiges Mittel, um Inhalte übersichtlich darzustellen oder Daten zu integrieren. Sie sind zentriert und vom Text abgesetzt einzubinden. Von umflossenen Tabellen ist abzusehen.

Wie auch Abbildungen sind Tabellen mit Beschriftungen zu versehen. Die Tabellenbeschriftung ist oberhalb der Tabelle anzuordnen, s. z. B. Tabelle \ref{tab:Beispieltab-1_DE}. Sie wird zentriert und im Blocksatz mit hängendem Einzug gesetzt, so dass die eigentliche Beschriftung linksbündig nach der Nummerierung beginnt, s. auch die Regelungen zu Abbildungen in Abschnitt \ref{cha:Images_DE}. 

\begin{table}[!htbp]
	\centering
	\caption[Zweispaltige Beispieltabelle mit unterschiedlichen Spaltenausrichtungen.]{Zweispaltige Beispieltabelle mit unterschiedlichen Spaltenausrichtungen (linke Spalte rechtsbündig, rechte Spalte linksbündig). Die Beschriftung hat einen hängenden Einzug.}
	\begin{tabular}{rl}
		\toprule
		\textbf{Arbeitsschritt} & \textbf{Teilarbeitsschritt} \\
		\midrule
		Einfacher Arbeitsschritt & Irgendetwas Einfaches \\
		& Noch irgendetwas Einfaches \\
		&\\
		Etwas komplizierterer Arbeitsschritt & Vorbereitende Arbeiten \\
		& Irgendetwas Kompliziertes \\
		& Noch irgendetwas Kompliziertes \\
		& Abschließende Arbeiten, die noch komplizierter sind \\
		\bottomrule
	\end{tabular}%
	\label{tab:Beispieltab-1_DE}%
\end{table}%

Die Beschriftung muss den Tabelleninhalt eindeutig beschreiben, so dass die Tabelle auch ohne Fließtext für sich stehend verständlich ist. Eventuell verwendete Symbole (Formelzeichen) sind in der Bildbeschriftung einzuführen. Bei langen Tabellenbeschriftungen gibt \LaTeX{} wie auch bei Abbildungen die Möglichkeit, eine Kurzform für das Tabellenverzeichnis zu definieren.

\begin{addmargin}[1cm]{1cm}
	\textit{\textbf{Hinweis:} Auf jede Tabelle muss im Fließtext Bezug genommen werden. Nicht referenzierte Tabellen sind zu entfernen.}
\end{addmargin}

Wird eine Tabelle aus einer anderen Quelle übernommen, ist diese in der Tabellenbeschriftung zu referenzieren. Es ist allerdings zu vermeiden, Tabellen aus anderen Quellen unverändert zu übernehmen, da sonst ein einheitlicher Tabellenstil nicht möglich ist. Besser ist es, nur den Inhalt einer Tabelle aus einer Quelle zu übernehmen und diesen in eine selbst gestaltete Tabelle zu integrieren. Auch in diesem Fall ist die Quelle zu referenzieren, z. B. in der Form " `Inhalte aus [XX]" ' (wobei \textit{XX} ein Platzhalter für die Referenznummer ist).

Tabellen dürfen nicht über die Textbreite hinausragen. In breiteren Tabellen kann in Grenzen die Schriftgröße verkleinert werden, s. Tabelle \ref{tab:Beispieltab-2_DE}. Es ist aber in jedem Fall darauf zu achten, dass die Lesbarkeit gewährleistet ist. Die Schriftgröße in Tabelle \ref{tab:Beispieltab-2_DE} ist grenzwertig. Kleiner sollte die Schriftgröße nicht werden. Wenn die Tabellen sehr breit sind, sollten sie im Querformat eingebunden werden.

\begin{table}[t]%[!htbp]
	\tiny
	\centering
	\caption[Siebenspaltige Beispieltabelle mit kleiner Schrift.]{Siebenspaltige Beispieltabelle mit kleiner Schrift, damit die Tabelle in die Textbreite passt. Kleiner als in diesem Beispiel sollte die Schriftgröße innerhalb einer Tabelle nicht sein.}
	\begin{tabular}{ccccccc}
		\toprule
		\textbf{Überschrift 1} & \textbf{Überschrift 2} & \textbf{Überschrift 3} & \textbf{Überschrift 4} & \textbf{Überschrift 5} & \textbf{Überschrift 6} & \textbf{Überschrift 7} \\
		\midrule
		Inhalt zu Überschrift 1 & Inhalt zu Überschrift 2 & Inhalt zu Überschrift 3 & Inhalt zu Überschrift 4 & ... & ... & ... \\
		Inhalt zu Überschrift 1 & Inhalt zu Überschrift 2 & Inhalt zu Überschrift 3 & Inhalt zu Überschrift 4 & ... & ... & ... \\
		Inhalt zu Überschrift 1 & Inhalt zu Überschrift 2 & Inhalt zu Überschrift 3 & ... & ... & ... & ... \\
		Inhalt zu Überschrift 1 & Inhalt zu Überschrift 2 & ... & ... & ... & ... & ... \\
		Inhalt zu Überschrift 1 & ... & ... & ... & ... & ... & ... \\
		... & ... & ... & ... & ... & ... & ... \\
		\bottomrule
	\end{tabular}%
	\label{tab:Beispieltab-2_DE}%
\end{table}%

Es gibt unzählige Möglichkeiten, wie eine Tabelle im Detail gestaltet werden kann. Für ein einheitliches Layout von Tabellen sollte auf vertikale Linien verzichtet werden, wenn sie nicht zwingend erforderlich sind. Oberhalb und unterhalb der Spaltenbeschriftungen und unterhalb der letzten Tabellenzeile ist eine horizontale Linie anzuordnen, wobei die oberste und unterste Linie etwas dicker sein soll, s. Tabelle \ref{tab:Beispieltab-1_DE} und Tabelle \ref{tab:Beispieltab-2_DE}. Zur vertikalen Abgrenzung können Leerzeilen eingefügt werden, s. Tabelle \ref{tab:Beispieltab-1_DE}. Alternativ können zusätzliche dünne horizontale Linien zur Unterteilung von Tabellenblöcken sinnvoll sein, s. Tabelle \ref{tab:Beispieltab-3_DE} und Tabelle \ref{tab:Beispieltab-4_DE}.

Die horizonalte Ausrichtung von Spalten (linksbündig, rechtsbündig oder zentriert) ist im Einzelfall abzuwägen. Die Tabellen \ref{tab:Beispieltab-1_DE} bis \ref{tab:Beispieltab-4_DE} zeigen die verschiedenen Möglichkeiten, auch in Kombination. Mitunter kann es zweckmäßig sein, Zeilen teilweise vertikal zu verbinden. Ein Beispiel dafür ist in Tabelle \ref{tab:Beispieltab-4_DE} dargestellt. Dort sind in den ersten beiden Spalten die Zellen vertikal verbunden, so dass in der dritten Spalte eine Zeilenaufteilung erkennbar ist. Die dünnen horizontalen Linien erhöhen in diesem Fall die Lesbarkeit der Tabelle.

\begin{table}[!hbpt]
	\tiny
	\centering
	\caption[Sechsspaltige Beispieltabelle mit kleiner Schrift.]{Sechsspaltige Beispieltabelle. Die dritte Spalte ist linksbündig, alle anderen Spalten sind zentriert.}
	\begin{tabular}{cclccc}
		\toprule
		\textbf{Zu\-ordnung} & \textbf{Pos.} & \textbf{Positionsbeschreibung} & \multicolumn{1}{c}{\textbf{GP [\EUR{}]}} & \multicolumn{1}{c}{\textbf{Prozent von GP}} & \multicolumn{1}{c}{\textbf{GP [\EUR{}]}} \\
		\midrule
		\multicolumn{1}{c}{\multirow{14}[0]{*}{A - Pos.}} & 3.5.10 & \multicolumn{1}{l}{Stb-Wand, C20/25; h = bis 19,05 m; d = 0,2 m herstellen} & \EUR{547.720,07} & \multicolumn{1}{c}{15,36\%} & \EUR{547.720,07} \\
		\multicolumn{1}{c}{} & 1.2.10 & \multicolumn{1}{l}{Vorhalten der Baustelle} & \EUR{420.308,58} & \multicolumn{1}{c}{27,15\%} & \EUR{968.028,65} \\
		\multicolumn{1}{c}{} & 4.2.10 & \multicolumn{1}{l}{Filigrandecke d = 0,05 m liefern und verlegen} & \EUR{312.057,83} & \multicolumn{1}{c}{35,90\%} & \EUR{1.280.086,48} \\
		\multicolumn{1}{c}{} & 4.8.20 & \multicolumn{1}{l}{Stb-Stützen, C30/37, b/d = 0,6/0,9, h = 18,75 m herstellen} & \EUR{226.309,35} & \multicolumn{1}{c}{42,25\%} & \EUR{1.506.395,83} \\
		\multicolumn{1}{c}{} & 4.2.20 & \multicolumn{1}{l}{Filigrandecke Maschinenhalle d = 0,05 m liefern u. verlegen} & \EUR{224.606,69} & \multicolumn{1}{c}{48,55\%} & \EUR{1.731.002,51} \\
		\multicolumn{1}{c}{} & 3.3.10 & \multicolumn{1}{l}{Bodenplatte, Stb C30/37, d = 25 cm herstellen} & \EUR{173.623,54} & \multicolumn{1}{c}{53,41\%} & \EUR{1.904.626,06} \\
		\multicolumn{1}{c}{} & 3.4.30 & \multicolumn{1}{l}{Stb-Stützen, C30/37, b/d = 0,6/2,8, h = 6,25 m herstellen} & \EUR{164.099,59} & \multicolumn{1}{c}{58,02\%} & \EUR{2.068.725,65} \\
		\multicolumn{1}{c}{} & 4.2.60 & \multicolumn{1}{l}{Aufbeton der Decke EG d = 0,35 m herstellen} & \EUR{151.737,34} & \multicolumn{1}{c}{62,27\%} & \EUR{2.220.462,99} \\
		\multicolumn{1}{c}{} & 2.1.20 & \multicolumn{1}{l}{Baugrubenaushub Bkl. 3-5, inkl. Abfuhr} & \EUR{134.779,63} & \multicolumn{1}{c}{66,05\%} & \EUR{2.355.242,63} \\
		\multicolumn{1}{c}{} & 3.1.20 & \multicolumn{1}{l}{Einzelfundamente Stb C20/25, b/l = 2,10/2,80; h = 1,50 m} & \EUR{126.460,84} & \multicolumn{1}{c}{69,60\%} & \EUR{2.481.703,47} \\
		\multicolumn{1}{c}{} & 4.6.10 & \multicolumn{1}{l}{Entwässerungsrinne: Breite 60 cm; Hmax = 1,50 m} & \EUR{115.367,81} & \multicolumn{1}{c}{72,83\%} & \EUR{2.597.071,28} \\
		\multicolumn{1}{c}{} & 3.2.10 & \multicolumn{1}{l}{Streifenfundamente Stb C20/25 b = 1,40 m; l = 10 m, 70 m; h = 1,50 m} & \EUR{91.621,95} & \multicolumn{1}{c}{75,40\%} & \EUR{2.688.693,23} \\
		\multicolumn{1}{c}{} & 1.1.10 & \multicolumn{1}{l}{Einrichten der Baustelle} & \EUR{88.494,12} & \multicolumn{1}{c}{77,88\%} & \EUR{2.777.187,35} \\
		\multicolumn{1}{c}{} & 4.2.50 & \multicolumn{1}{l}{Aufbeton Maschinenhalle Stb C20/25 d = 0,15 m herstellen} & \EUR{86.053,13} & \multicolumn{1}{c}{80,30\%} & \EUR{2.863.240,48} \\
		\midrule
		\multicolumn{1}{c}{\multirow{9}[0]{*}{B - Pos.}} & 4.7.20 & \multicolumn{1}{l}{Stb-Stützen, C30/37, b/d = 0,6/0,6, h = 18,75 m herstellen} & \EUR{83.945,37} & \multicolumn{1}{c}{82,65\%} & \EUR{2.947.185,85} \\
		\multicolumn{1}{c}{} & 4.4.10 & \multicolumn{1}{l}{Deckenunterzüge, Stb C35/45 , b/d = 0,6/0,8 m herstellen} & \EUR{72.332,54} & \multicolumn{1}{c}{84,68\%} & \EUR{3.019.518,39} \\
		\multicolumn{1}{c}{} & 1.3.10 & \multicolumn{1}{l}{Räumen der Baustelle} & \EUR{71.922,97} & \multicolumn{1}{c}{86,70\%} & \EUR{3.091.441,36} \\
		\multicolumn{1}{c}{} & 4.3.10 & \multicolumn{1}{l}{Spannbetonbinder C35/45 Hallendecke herstellen} & \EUR{65.125,35} & \multicolumn{1}{c}{88,52\%} & \EUR{3.156.566,71} \\
		\multicolumn{1}{c}{} & 3.1.10 & \multicolumn{1}{l}{Einzelfundamente Stb C20/25, b/l = 1,60 /2,80; h = 1,50 m} & \EUR{52.688,79} & \multicolumn{1}{c}{90,00\%} & \EUR{3.209.255,50} \\
		\multicolumn{1}{c}{} & 3.2.20 & \multicolumn{1}{l}{Streifenfundamente Stb C20/25, b/l = 0,6/4,8; h = 1,50 m} & \EUR{46.704,40} & \multicolumn{1}{c}{91,31\%} & \EUR{3.255.959,90} \\
		\multicolumn{1}{c}{} & 4.7.10 & \multicolumn{1}{l}{Stb-Stützen, C30/37, b/d = 0,6/0,6, h = 10,95 m herstellen} & \EUR{44.146,46} & \multicolumn{1}{c}{92,55\%} & \EUR{3.300.106,36} \\
		\multicolumn{1}{c}{} & 2.2.20 & \multicolumn{1}{l}{Trag.- und Frostschutzschicht liefern und herstellen} & \EUR{43.487,87} & \multicolumn{1}{c}{93,77\%} & \EUR{3.343.594,22} \\
		\multicolumn{1}{c}{} & 4.8.10 & \multicolumn{1}{l}{Stb-Stützen, C30/37, b/d = 0,6/0,9, h = 6,25 m herstellen} & \EUR{42.400,87} & \multicolumn{1}{c}{94,96\%} & \EUR{3.385.995,09} \\
		\bottomrule
	\end{tabular}%
	\label{tab:Beispieltab-3_DE}%
\end{table}%

\begin{table}[htbp!]
	\centering
	\caption[Dreispaltige Beispieltabelle mit Gruppenstruktur.]{Dreispaltige Beispieltabelle mit Gruppenstruktur.}
	\begin{tabularx}{\textwidth}{l l X}
		\toprule
		\textbf{Fachbereich} & \textbf{Kernaufgabe} & \textbf{Aufgabenbereiche}\\
		\midrule
		Project Development & Entwicklung & Projektentwicklung, Site Assessment, Feasibility Studies\\
		\midrule
		Project Management Coordination & Projektleitung & Ansprechpartner gegenüber dem Kunden, Veranlassung der Rechnungsstellung, Projektablaufplanung, Umsetzung der Vertragspflichten, ggf.\ Disposition der Baustelleneinrichtung, ggf.\ Krandisposition\\
		\midrule
		Civil Works & Erdarbeiten & Entwurf, Kontrolle und Abnahme des Fundamentdesigns und des Fundamentbaus; Entwurf, Kontrolle und Abnahme der Zuwegung und der Kranstellfläche\\
		\bottomrule
	\end{tabularx}
	\label{tab:Beispieltab-4_DE}
\end{table}


Tabelle \ref{tab:Beispieltab-5_DE} stellt ein Beispiel für eine mehrseitige Tabelle dar, die mithilfe der 'longtable'-Latex-enviroment erstellt wurde.


\begin{longtable}{lcc}
	\caption{Windenergieproduktion nach Ländern \cite{elec}}\\
	\toprule
	\textbf{Land}             & \textbf{Windenergieerzeugung} & \textbf{Anteil von Gesamtenergie} \\
	& \textbf{(TWh)}                & \textbf{(\%)}                     \\
	\midrule
	\endfirsthead
	\caption{(Fortsetzung) Windenergieproduktion nach Ländern \cite{elec}}\\
	\toprule
	\textbf{Land}             & \textbf{Windenergieerzeugung} & \textbf{Anteil von Gesamtenergie} \\
	& \textbf{(TWh)}                & \textbf{(\%)}                     \\
	\midrule
	\endhead
	World                     & 2098.46                       & 7.3                               \\
	China                     & 762.70                        & 8.6                               \\
	United States             & 434.30                        & 10.1                              \\
	Germany                   & 124.91                        & 22.3                              \\
	United Kingdom            & 80.26                         & 24.6                              \\
	Brazil                    & 80.00                         & 11.8                              \\
	India                     & 70.05                         & 3.8                               \\
	Spain                     & 61.90                         & 21.7                              \\
	France                    & 38.56                         & 8.2                               \\
	Canada                    & 38.13                         & 6.0                               \\
	Turkey                    & 35.14                         & 10.8                              \\
	Sweden                    & 33.33                         & 19.4                              \\
	Australia                 & 31.65                         & 11.6                              \\
	Netherlands               & 21.49                         & 17.7                              \\
	Italy                     & 20.28                         & 7.1                               \\
	Mexico                    & 20.23                         & 5.8                               \\
	Poland                    & 19.73                         & 11.0                              \\
	Denmark                   & 18.93                         & 55.0                              \\
	Norway                    & 15.90                         & 10.4                              \\
	Argentina                 & 14.02                         & 9.8                               \\
	Portugal                  & 13.28                         & 28.3                              \\
	Belgium                   & 12.16                         & 12.8                              \\
	Finland                   & 11.72                         & 16.1                              \\
	Ireland                   & 11.18                         & 33.3                              \\
	Greece                    & 10.87                         & 20.8                              \\
	South Africa              & 9.69                          & 4.5                               \\
	Chile                     & 8.82                          & 10.6                              \\
	Japan                     & 8.17                          & 0.8                               \\
	Vietnam                   & 8.04                          & 3.1                               \\
	Romania                   & 6.97                          & 12.6                              \\
	Austria                   & 6.96                          & 10.6                              \\
	Russia                    & 5.51                          & 0.5                               \\
	Morocco                   & 5.11                          & 12.4                              \\
	Uruguay                   & 4.80                          & 30.5                              \\
	Egypt                     & 4.23                          & 2.0                               \\
	Pakistan                  & 4.12                          & 2.7                               \\
	Taiwan                    & 3.50                          & 1.2                               \\
	South Korea               & 3.37                          & 0.6                               \\
	New Zealand               & 3.20                          & 7.3                               \\
	Thailand                  & 3.00                          & 1.6                               \\
	Kazakhstan                & 2.28                          & 2.0                               \\
	Croatia                   & 2.28                          & 16.1                              \\
	Peru                      & 1.96                          & 3.3                               \\
	Kenya                     & 1.68                          & 13.6                              \\
	Lithuania                 & 1.59                          & 38.0                              \\
	Jordan (2021)             & 1.57                          & 7.2                               \\
	Bulgaria                  & 1.44                          & 2.9                               \\
	Costa Rica                & 1.36                          & 10.8                              \\
	Dominican Republic (2021) & 1.29                          & 7.3                               \\
	Iran                      & 1.25                          & 0.4                               \\
	Philippines               & 0.95                          & 0.8                               \\
	Serbia                    & 0.94                          & 2.7                               \\
	Ukraine                   & 0.83                          & 0.7                               \\
	Honduras (2021)           & 0.80                          & 6.7                               \\
	Panama (2021)             & 0.66                          & 5.9                               \\
	Czech Republic            & 0.65                          & 0.8                               \\
	Estonia                   & 0.63                          & 7.5                               \\
	Nicaragua (2021)          & 0.63                          & 13.6                              \\
	Ethiopia (2021)           & 0.61                          & 4.2                               \\
	Hungary                   & 0.61                          & 1.7                               \\
	Bolivia                   & 0.48                          & 4.3                               \\
	Tunisia                   & 0.42                          & 2.1                               \\
	Mongolia                  & 0.42                          & 5.8                               \\
	Bosnia and Herzegovina    & 0.39                          & 2.4                               \\
	Sri Lanka (2021)          & 0.37                          & 2.3                               \\
	Indonesia                 & 0.36                          & 0.1                               \\
	Guatemala (2021)          & 0.32                          & 2.2                               \\
	Luxembourg                & 0.31                          & 25.2                              \\
	Montenegro                & 0.31                          & 9.4                               \\
	Jamaica (2021)            & 0.30                          & 6.9                               \\
	Senegal (2021)            & 0.25                          & 4.5                               \\
	Cyprus                    & 0.22                          & 4.2                               \\
	Latvia                    & 0.19                          & 3.8                               \\
	\bottomrule
	\label{tab:Beispieltab-5_DE}
\end{longtable}

